\section{Bounds on the radius and occupation times}\label{sec:coarse}
We first bound the range and local times at their natural scales.

\begin{proposition}\label{prop:coarse}
Let $d\geq2$ be an integer, let $0<\eps<1/d$, and let $(X_j)_{j\geq0}$ be centrally
excited random walk on $\Z^d$, started at the origin, with first-departure probabilities
$q_x(\pm e_i)=1/(2d)\mp\eps u_x^i/2$ and simple random walk steps
on subsequent departures and at the origin. For every $p>0$, there are constants
$c,C>0$, depending only on $d,\eps,p$, such that for every integer
$n\geq2$, with probability at least $1-Cn^{-p}$,
\begin{equation}\label{eq:coarse}
cN^d\leq R_n\leq CN^d ,\qquad
cN\leq M_n\leq CN ,\qquad cN\leq H_n\leq CN\, ,
\end{equation}
where $N=n^{1/(d+1)}$, $R_n$ counts distinct departure sites,
$M_n$ is the maximum departure local time, and
$H_n=\max_{j\leq n}|X_j|$.
\end{proposition}

\subsection{A radial test}
The discrete Laplacian of the truncated potential is supported near
the truncation level.
Throughout this section, $F$ is the tail in \eqref{eq:Fdef} for $D_n$.
Fix $b_d=\lceil\sqrt d/2\rceil+6$, and let
\begin{equation}\label{eq:shell-def}
M_{\rm sh}(r)=\max_{x:||x|-r|\leq3}\ell_n(x)\, .
\end{equation}
The potential kernel $b$ was defined in Section~\ref{sec:local}.
For $r>0$, write
\begin{equation}\label{eq:radial-test}
\phi_r(x)=(b(x)-h_d(r))_+ ,\qquad
h_d(r)=\begin{cases}
\dfrac2\pi\log r+\kappa,&d=2,\\[3pt]
-\dfrac2{(d-2)\omega_d}r^{2-d},&d\geq3
\end{cases}\, .
\end{equation}

\begin{lemma}\label{lem:radial}
Let $d\geq2$ be an integer, let $0<\eps<1/d$, and let $(X_j)_{j\geq0}$ be centrally
excited random walk started at the origin, with first-departure probabilities
$q_x(\pm e_i)=1/(2d)\mp\eps u_x^i/2$, with simple symmetric
steps on later departures and at the origin.
Let $b_d=\lceil\sqrt d/2\rceil+6$, and let
$F(s)=\sigma_d^{-1}\int_{D_n\cap\{|v|>s\}}|v|^{1-d}\dd v$.
For every $p>0$, there are $r_0=r_0(d)>2b_d$ and $C=C(d,\eps,p)>0$
such that, for every integer $n\geq2$, with probability at least
$1-Cn^{-p}$, simultaneously for
all integers $r_0\leq r\leq n$,
\begin{equation}\label{eq:radial}
\begin{split}
F(r+b_d)\leq{}&CM_{\rm sh}(r)
 [F(r-b_d)-F(r+b_d)]\\
&+C\left[\sqrt{M_nr^{1-d}F(r-b_d)L}+r^{1-d}L\right]\, .
\end{split}
\end{equation}
Here $L=\log(n+2)$ and
$M_{\rm sh}(r)=\max_{x:\,||x|-r|\leq3}\ell_n(x)$.
\end{lemma}

\begin{proof}
The asymptotics \eqref{eq:kernel-asymptotics} imply that, for all
sufficiently large $r$,
\begin{equation}\label{eq:levelsets}
\phi_r(x)=0\quad\hbox{if }|x|\leq r-1 ,\qquad
\phi_r(x)=b(x)-h_d(r)\quad\hbox{if }|x|\geq r+1\, .
\end{equation}
Indeed, the radial leading term increases by order $r^{1-d}$
over a unit interval, while its error is $O(r^{-d})$.
The radial upper and lower envelopes are increasing for large radii.
The remaining finite set of inner sites is covered by $h_2(r)\to\infty$
in the plane, and by $h_d(r)\to0$ and $b(x)<0$ in higher dimensions.

Every increment of $\phi_r$ vanishes from $|x|\leq r-2$.
For $|x|>r+3$, exact harmonicity and \eqref{eq:gradient} give
\begin{equation}\label{eq:radial-drift}
(P-I)\phi_r(x)=0 ,\qquad
u_x\cdot D\phi_r(x)\geq c_d|x|^{1-d}\, .
\end{equation}
On the shell $||x|-r|\leq3$, both $|(P-I)\phi_r(x)|$ and
$|D\phi_r(x)|$ are at most $C_dr^{1-d}$, because the positive-part
map is $1$-Lipschitz.

The Dynkin martingale $W^{(r)}$ for $\phi_r(X_j)$ has increments
bounded by $C_dr^{1-d}$ and bracket
\begin{equation}\label{eq:radialbracket}
\langle W^{(r)}\rangle_n
\leq C_d\sum_{|x|>r-2}\ell_n(x)|x|^{2-2d}
\leq C_dM_nr^{1-d}F(r-b_d)\, .
\end{equation}
For the second inequality, each departure cell in the sum lies outside
$r-b_d$, its weighted integral is comparable to $|x|^{1-d}$, and
$|x|^{2-2d}\leq C_dr^{1-d}|x|^{1-d}$.
The martingale $r^{d-1}W^{(r)}$ has bounded increments and bracket
at most $C_dn$. Dyadic variance bounds, including the bound below one,
and a union bound over integer $r$ give
\begin{equation}\label{eq:radialmart}
|W_n^{(r)}|\leq C\left[
\sqrt{M_nr^{1-d}F(r-b_d)L}+r^{1-d}L\right]\, .
\end{equation}

Since $\phi_r(0)=0$ and $\phi_r(X_n)\geq0$, Dynkin's formula gives
\begin{equation}\label{eq:radial-source}
c_d\eps\sum_{\substack{x\in A_n\\|x|>r+3}}|x|^{1-d}
\leq C_dr^{1-d}\sum_{||x|-r|\leq3}\ell_n(x)+|W_n^{(r)}|\, .
\end{equation}
The first-departure contribution on the shell is absorbed by its local
time. The left side bounds a positive constant times $F(r+b_d)$.
Every shell cell lies between $r-b_d$ and $r+b_d$, so their number
is at most
\begin{equation}\label{eq:shell-count}
\sigma_d(r+b_d)^{d-1}[F(r-b_d)-F(r+b_d)]\, .
\end{equation}
Multiplying by $r^{1-d}$ proves \eqref{eq:radial}.
\end{proof}

\subsection{Reducing the exterior volume}
The spherical average and H\"older continuity bound the largest local
time on a shell in terms of the exterior volume.
Write $s_*=R_n^{1/d}$ and $\delta=Ce_n(M_n)$.
Suppose for now that
\begin{equation}\label{eq:preconditions}
s_*\geq cN ,\qquad M_n\leq Cs_* ,\qquad
\delta\leq C\sqrt{s_*}L\, .
\end{equation}
We verify these estimates before using the argument.
Let
\begin{equation}\label{eq:alpha-beta}
\alpha=\frac1{2d-1} ,\qquad \beta=\frac{d-1}{2d-1}\, .
\end{equation}
For a fixed $B_0\geq2$ and $r\in[s_*,B_0s_*]$, we claim that
\begin{equation}\label{eq:shell}
M_{\rm sh}(r)\leq CB_0^\beta s_*^{1-\alpha}
 [F(r-b_d)+\delta]^\alpha\, ,
\end{equation}
with $C$ independent of $B_0$.
On the event of Lemma~\ref{lem:local}, $U_{D_n}+\delta\geq0$
throughout the approximation region.
Let $h$ be its maximum on a sphere with
$s_*/2\leq t\leq B_0s_*+3$. For fixed $B_0$, these spheres lie
inside $B(0,2n)$ for large $n$.
By \eqref{eq:holder}, the function is at least $h/2$ on a spherical
cap of chord radius $ch^2/s_*$. The constant $c$ can be chosen so
that this radius is at most $t/2$, since $h\leq Cs_*$.
The cap area is at least $c_d(h^2/s_*)^{d-1}$. Hence
\begin{equation}\label{eq:cap-average}
2d\eps F(t)+\delta\geq
\frac{ch^{2d-1}}{s_*^{d-1}t^{d-1}}\, .
\end{equation}
Apply this bound on the sphere through each shell site and use the
monotonicity of $F$. This proves \eqref{eq:shell}.

We next reduce $F$ to the threshold
\begin{equation}\label{eq:floor}
\Lambda=Cs_*^{2-d}L\, .
\end{equation}
Use a grid of spacing $2b_d$, with integer $r$ at each midpoint.
For $r\geq s_*$, the error in \eqref{eq:radial} is at most
\begin{equation}\label{eq:halving-error}
C\sqrt{s_*^{2-d}F(r-b_d)L}+Cs_*^{1-d}L\, .
\end{equation}
If $F(r-b_d)\leq A$ and $A\geq\Lambda$, then the error is at most
$A/4$, after increasing the constant in \eqref{eq:floor}.
Unless $F(r+b_d)\leq A/2$, the decrease is at least $cA/K(A)$,
where
\begin{equation}\label{eq:halving-K}
K(A)=CB_0^\beta s_*^{1-\alpha}(A+\delta)^\alpha+1\, .
\end{equation}
A halving therefore takes at most $CK(A)$ grid steps.
For an initial bound $A_0\leq Cs_*$, the number of dyadic levels
above $\Lambda$ is $O_d(L)$, because $s_*\leq n^{1/d}$.
The sum of their $A^\alpha$ contributions is bounded by
$CA_0^\alpha$. Thus the radial distance required to reach $\Lambda$
is at most
\begin{equation}\label{eq:halving-cost}
CB_0^\beta s_*^{1-\alpha}A_0^\alpha
+C_{B_0}s_*^{1-\alpha/2}L^{1+\alpha}+CL\, .
\end{equation}
The first constant is independent of $B_0$.
The last two terms are $o(s_*)$, uniformly under \eqref{eq:preconditions}.
This calculation is valid whenever the stated distance fits inside
the window $[s_*,B_0s_*]$.

\subsection{Crossings through a small set}
An interval with few distinct departure sites cannot make a long
outward crossing.
The compensated position
\begin{equation}\label{eq:vector-def}
Z_t=X_t+\eps\sum_{j<t}I_ju_{X_j}\, .
\end{equation}
is a vector martingale with bounded increments.
Azuma's inequality in each coordinate and a union bound over deterministic
intervals give, with probability at least $1-Cn^{-p}$,
\begin{equation}\label{eq:vector}
|Z_t-Z_s|\leq C\sqrt{(t-s)L}
\qquad(0\leq s<t\leq n)\, .
\end{equation}
Consequently every projection chosen from the path is also controlled.

Consider an interval $[s,t]$ whose departure sites satisfy
$v\cdot X_j>b>0$ and $|X_j|<r$, where $|v|=1$ and $b/r\geq\kappa>0$.
Suppose that $v\cdot(X_t-X_s)\geq h>0$.
Let $k$ be the number of fresh departures and $m$ the number of
distinct departure sites on the interval.
Equation~\eqref{eq:interval} bounds its duration by
$m(C\eps k^{1/d}+C\lambda_n)$.
Each fresh departure has projected inward mean at least $\kappa\eps$.
By \eqref{eq:vector},
\begin{equation}\label{eq:crossing-before-young}
(h+\kappa\eps k)^2\leq Cm(\eps k^{1/d}+\lambda_n)L\, .
\end{equation}
Young's inequality with powers $2d$ and $2d/(2d-1)$ absorbs half
the $(\kappa\eps k)^2$ term and gives
\begin{equation}\label{eq:crossing}
h^2\leq C_\kappa\bigl(m^\gamma L^\gamma+m\lambda_nL\bigr) ,
\qquad\gamma=\frac{2d}{2d-1}<2\, .
\end{equation}
The event on which this implication is used is simultaneous over
all deterministic intervals and does not depend on the later choice
of $v,b,r$.

\subsection{Proof of the occupation and radius bounds}
We now close the estimates without a prior bound on the outer radius.

\begin{proof}[Proof of Proposition~\ref{prop:coarse}]
Intersect the events in Lemmas~\ref{lem:local} and~\ref{lem:radial}
with the event in \eqref{eq:vector}.
Since $n\leq M_nR_n$, equation~\eqref{eq:M} gives
\begin{equation}\label{eq:sstar-lower}
n\leq C\eps s_*^{d+1}+C\lambda_ns_*^d\, .
\end{equation}
This forces $s_*\geq cN$ for large $n$: if $c$ is sufficiently small,
then the first term for $s_*<cN$ is at most $n/2$ and the second is $o(n)$.
It follows that $\lambda_n=o(s_*)$, $M_n\leq Cs_*$, and
$\delta\leq C\sqrt{s_*}L$. This verifies \eqref{eq:preconditions}.

Start the halving argument at $\lceil s_*\rceil$, where
$F\leq s_*/\sigma_d$.
Choose $B_0\geq8$ so large that the leading term
$CB_0^\beta s_*$ in \eqref{eq:halving-cost} is less than $B_0s_*/4$.
This is possible because $\beta<1$.
For sufficiently large $n$, the remaining terms are also less than
$B_0s_*/4$. The halvings therefore finish before $(B_0-2)s_*$, and
\begin{equation}\label{eq:coarse-tail}
F((B_0-2)s_*)\leq Cs_*^{2-d}L\, .
\end{equation}
All integer radii used lie below $n$ for large $n$, since
$s_*\leq n^{1/d}$ and $B_0$ is fixed.

Suppose for a contradiction that $H_n\geq(B_0+1)s_*$.
Let $\tau$ be the first time the path reaches or exceeds this radius, let
$v=X_\tau/|X_\tau|$, and let $b=(B_0-1)s_*$.
Choose the last entrance into the half-space $v\cdot x>b$ before
$\tau$:
\begin{equation}\label{eq:coarse-entrance}
s=1+\max\{j<\tau:v\cdot X_j\leq b\}\, .
\end{equation}
The set contains zero. At entrance, $v\cdot X_s\leq b+1$.
For $s\leq j<\tau$, the departure site $X_j$ has projection greater
than $b$ and radius less than $(B_0+1)s_*$. The net projected displacement
is at least $s_*$ for large $n$.
Their $m$ distinct cells belong to $D_n$, including when $\tau=n$.
By \eqref{eq:Fmass} and \eqref{eq:coarse-tail},
\begin{equation}\label{eq:coarse-cellcount}
m\leq\sigma_d((B_0+1)s_*+\sqrt d/2)^{d-1}
F(b-\sqrt d/2)\leq C_{B_0}s_*L\, .
\end{equation}
The crossing bound \eqref{eq:crossing} would imply
\begin{equation}\label{eq:coarse-contradiction}
s_*^2\leq C_{B_0}\bigl(s_*^\gamma L^{2\gamma}+s_*L^4\bigr)
=o(s_*^2)\, .
\end{equation}
We conclude that
\begin{equation}\label{eq:HvsR}
H_n\leq CR_n^{1/d}\, .
\end{equation}

It remains to bound $R_n$ in terms of $n$.
Let $S=H_n+\sqrt d$, so that $D_n\subset B(0,S)$.
Integrating the spherical average in \eqref{eq:newton} gives
\begin{equation}\label{eq:massidentity}
\int_{B(0,S)}U_{D_n}(y)\dd y
=2d\eps\sigma_d\int_0^S r^{d-1}F(r)\dd r
=2\eps\int_{D_n}|v|\dd v\, .
\end{equation}
Since $\int\widetilde\ell_n=n$, the approximation error gives
\begin{equation}\label{eq:coarse-masserror}
\left|n-2\eps\int_{D_n}|v|\dd v\right|
\leq\omega_dS^d\delta\leq Cs_*^{d+1/2}L=o(s_*^{d+1})\, .
\end{equation}
Among sets of volume $R_n$, the centered ball minimizes the integral
of $|v|$. Therefore
\begin{equation}\label{eq:radial-packing}
\int_{D_n}|v|\dd v\geq
\frac d{d+1}\omega_d^{-1/d}R_n^{1+1/d}\, .
\end{equation}
Absorb the error in \eqref{eq:coarse-masserror} to obtain $R_n\leq CN^d$.
Equations~\eqref{eq:M} and~\eqref{eq:HvsR} imply $M_n,H_n\leq CN$.
The lower bounds follow from $R_n\geq cN^d$,
$R_n\leq C_d(H_n+1)^d$, and $n\leq M_nR_n$.
Increasing constants covers bounded $n$.
\end{proof}
